\documentclass{situationist-preprint}

\title{Parfit on What Matters}
\msid{PID-03-b}
\msfamily{PID-03}
\msdiscipline{Personal Identity and the Self}
\mstype{Companion manuscript}
\mskeywords{Parfit; fission; what matters; Relation R; reductionism; connectedness}

\begin{document}
\maketitle

\begin{abstract}
PID-03-a argues from fission that identity is not what matters in survival. The conclusion is Parfit's, and so is the case. We set out where the argument of PID-03-a agrees with Parfit's and where it departs from it. It agrees that fission produces two persons, neither identical with the original, each having what matters. It departs from Parfit in three respects: identity has a determinate answer in fission; identity keeps a small residual weight; and within what Parfit called Relation R, connectedness carries most of what matters, while continuity carries little.
\end{abstract}

\section{Agreement}

\citet[ch.~12]{Parfit1984} argued that in fission, where each hemisphere of a person's brain is transplanted into a body of its own, the original cannot be identical with both successors, since they are two; she has no better claim to be one than the other; and yet the outcome cannot be as bad as death, since if either transplant had occurred alone she would have survived. Since what happens to the original is about as good as ordinary survival, and does not contain identity, identity is not what matters. PID-03-a reaches the same conclusion by the same route: Lemma~3.1 (identity with neither), Theorem~3.3 (each successor has what matters) and Corollary~3.4 (identity is not what matters). The decisive premise, in both, is that whether a later being has what matters depends on its relation to the original and not on whether a rival exists (PID-03-a, Assumption~2.6), a premise Parfit expressed by saying that a double success cannot be a failure. \citet{Lewis1976} sought to keep identity and what matters together by counting two persons who share their stages before fission; PID-03-a, Remark~3.5, notes that what matters is unaffected on either count, which is the point on which Parfit and PID-03-a agree.

\section{Three departures}

\paragraph{Determinacy.} Parfit held that questions about identity in puzzle cases can be empty: once the facts about bodies, brains and psychological connections are known, there may be no further fact about whether the original survives, and it does not matter which description we choose \citep[ch.~11]{Parfit1984}. PID-03-a gives a determinate answer. The criterion of PID-01, psychological continuity with bodily continuity as a tie-breaker in branching, implies that in symmetric fission the original is neither successor. The answer is determinate because the criterion settles every case it covers, including those in which identity fails. What makes the conclusion about what matters forceful is not that the identity question is empty but that its answer, whatever it is, leaves the value of the outcome unaffected.

\paragraph{Residual weight.} Parfit's most radical statements say that identity is not what matters at all. PID-06-a reaches a weaker conclusion: identity carries at most a small weight $\varepsilon$ (Corollary~3.4). Fission shows that identity is not needed for nearly all of what matters, but not that its loss costs nothing. Some goods are goods of uniqueness: being the only one who stands in a certain relation to one's partner, children or commitments. A successor who shares these with another successor cannot have them in the same way. The residue is small, and it does not restore identity to the place of what matters.

\paragraph{Connectedness and continuity.} Parfit's Relation R is psychological connectedness and/or continuity, with any cause \citep[ch.~12]{Parfit1984}. He sometimes allowed that connectedness may matter more than continuity, and the long-chain case supports this. PID-06-a, Theorem~3.3, shows that if the contributions are approximately additive, connectedness carries more than identity and continuity together. Continuity without connectedness, as between a person and her distant past self, carries little. This sharpens Parfit's view rather than abandoning it, and it is the version used in PID-07 and PID-14.

\section{Conclusion}

PID-03-a is Parfit's argument, stated with a determinate criterion of identity and a graded account of what matters. Its conclusion, that identity is not what matters, is Parfit's, and does not depend on his claim that identity questions may be empty.

\begin{thebibliography}{99}\small
\bibitem[Parfit(1984)]{Parfit1984} Parfit, D. (1984). \emph{Reasons and Persons}. Oxford University Press.
\bibitem[Lewis(1976)]{Lewis1976} Lewis, D. (1976). Survival and identity. In A.~O. Rorty (ed.), \emph{The Identities of Persons} (pp.~17--40). University of California Press.
\end{thebibliography}

\end{document}
